\documentclass[11pt]{article}
\usepackage[margin=1in]{geometry}
\usepackage{amsmath, amssymb}
\usepackage{booktabs}
\usepackage{hyperref}
\usepackage{listings}
\usepackage{xcolor}
\lstset{basicstyle=\ttfamily\small, breaklines=true, columns=fullflexible,
        backgroundcolor=\color{black!5}, frame=single, numbers=none}

\title{Independent Replication of\\ ``On Quantum Algorithms for Noncommutative Hidden Subgroups''\\ \small (Ettinger \& H{\o}yer, arXiv:quant-ph/9807029, 1998)}
\author{QC-200 Replication Wave (rick-stevens-ai / OpenClaw subagent)}
\date{2026-07-05}

\begin{document}
\maketitle

\begin{abstract}
We independently replicate the central quantum algorithm of Ettinger \& H{\o}yer
(1998) for the dihedral Hidden Subgroup Problem (HSP), using real Qiskit
statevector simulation for $D_N$ with $N\in\{4,8,16,32\}$. The paper's Lemma~4
measurement distribution is reproduced to machine precision
($\mathrm{TV}=10^{-16}$--$10^{-15}$) across all sizes, and the
Ettinger--H{\o}yer post-processing recovers the hidden reflection parameter $k_0$
with observed query complexity well inside the paper's $O(\log N)$ bound. We
verify the classical-vs-quantum separation on the same instances (uniform
baseline never reaches $2/3$ success within the swept budget).
\textbf{Verdict: REPLICATED.} During the replication we identified one
paper-side inconsistency: the sin-based post-processing on $b{=}1$ outcomes
described in the proof of Theorem~3 has expectation zero under the paper's own
distribution and therefore cannot discriminate $k_0$; only the $b{=}0$/cos
branch is a valid estimator. We document this in \S\ref{sec:paper-issue}.
\end{abstract}

\section{Paper summary}
The paper by Ettinger \& H{\o}yer is the first known quantum algorithm for a
non-Abelian Hidden Subgroup Problem. It targets the dihedral group $D_N =
\mathbb{Z}_N \rtimes_\varphi \mathbb{Z}_2$ with multiplication
$(a_1,b_1)(a_2,b_2) = (a_1 + \varphi(b_1)(a_2), b_1 + b_2)$ where
$\varphi(1)(a) = -a$. Given a black-box $\gamma: D_N \to R$ promised to be
constant and distinct on left cosets of a hidden subgroup $H \leq D_N$, the
problem asks for a generating set of $H$.

The main results are:
\begin{itemize}
  \item \textbf{Theorem 2} reduces the general dihedral HSP to identifying
        reflection subgroups of the form $H = \{(0,0),(k_0,1)\}$ using
        $\Theta(\log N)$ oracle queries with success probability $1 - 2/N$.
  \item \textbf{Theorem 3} solves the reflection-subgroup case with at most
        $89 \log(N)+7$ oracle queries and success probability
        $1 - 1/(2N)$.
  \item The core quantum subroutine is $V^\gamma = (F_N \otimes W \otimes I)
        \cdot U_\gamma \cdot (F_N^{-1} \otimes W \otimes I)$ applied to
        $|0\rangle|0\rangle|0\rangle$; here $F_N$ is the QFT on
        $\mathbb{Z}_N$ and $W$ the Walsh--Hadamard on $\mathbb{Z}_2$.
  \item \textbf{Lemma 4} states the measurement distribution:
        $\Pr[(a,0)] = \tfrac{1}{N}\cos^2(\pi k_0 a / N)$ and
        $\Pr[(a,1)] = \tfrac{1}{N}\sin^2(\pi k_0 a / N)$.
  \item \textbf{Theorem 5} gives the post-processing: with
        $m \geq \lceil 64\ln N\rceil$ samples $\{z_i\}$ from
        the $b{=}0$ conditional distribution
        $\Pr[Z{=}z] = \alpha \cos^2(\pi k_0 z / N)$, the estimator
        $\tilde k = \arg\max_{k\in[1..N/2]} \sum_i \cos(2\pi k z_i / N)$
        satisfies $\tilde k = \min(k_0, N-k_0)$ with probability $\geq 1 - 1/(2N)$.
\end{itemize}
Post-processing runs in $O(N)$ time and does not admit a known
polynomial-in-$\log N$ algorithm --- the paper explicitly leaves that as an
open question. The paper's contribution is thus a \emph{query-complexity} result
(linear in $\log N$), not a total polynomial-time algorithm.

\section{Claims table}
\begin{center}
\begin{tabular}{@{}p{0.9cm} p{6.4cm} p{2.4cm} p{1.8cm} p{2.6cm}@{}}
\toprule
ID & Claim & Type & Testable? & Tested here? \\
\midrule
C1 & Lemma 4 measurement distribution
     $\Pr[(a,0)] = \tfrac{1}{N}\cos^2(\pi k_0 a/N)$ under $V^\gamma$
     & Analytical + simulable & Yes & \textbf{Yes} \\
C2 & Theorem 3: $\tilde k=\min(k_0,N-k_0)$ recovered with $\geq 1-1/(2N)$
     probability using $\leq 89\log N + 7$ oracle queries
     & Query-complexity & Yes (small $N$) & \textbf{Yes} \\
C3 & Classical black-box lower bound: no efficient classical algorithm
     & Complexity-theoretic & Not directly; simulable baseline & Partial
     (uniform baseline never reaches $2/3$ in swept budget) \\
C4 & Theorem 2 reduction (general dihedral $\to$ reflection case)
     via classical Abelian post-processing
     & Constructive & Yes but auxiliary & Not tested (out of scope of the reproducible core) \\
C5 & Post-processing is $O(N)$-time; no known $\mathrm{poly}(\log N)$ algorithm
     & Complexity remark & No (open problem) & N/A \\
\bottomrule
\end{tabular}
\end{center}

\section{Method (numbered)}
\begin{enumerate}
  \item Downloaded PDF from
        \url{https://arxiv.org/pdf/quant-ph/9807029} (162~kB), extracted text with
        \texttt{pdftotext -layout} for reading; skimmed for the reproducible core.
  \item Created a Python venv (\texttt{python3 -m venv .venv}) and installed
        \texttt{qiskit==2.5.0} and \texttt{numpy==2.5.1}. Free tools, no LLM inference required.
  \item Wrote \texttt{report/evidence/ettinger\_hoyer\_dihedral.py} implementing:
    \begin{enumerate}
      \item Oracle $U_\gamma$ constructed as an explicit unitary matrix
        on $n{+}1{+}(n{+}1)$ qubits ($n=\log_2 N$), keyed off a table of
        left-coset representatives of $H = \{(0,0),(k_0,1)\}$.
        Unitarity is asserted numerically ($UU^\dagger = I$).
      \item The exact DFT $F_N$ on $\mathbb{Z}_N$ as an explicit unitary
        (\emph{not} QFT-with-swaps), so the mathematical statement of Lemma~4
        is tested in its cleanest form.
      \item Full $V^\gamma = (F_N \otimes W \otimes I)\, U_\gamma\,
        (F_N^{-1} \otimes W \otimes I)$ applied to $|0\rangle|0\rangle|0\rangle$
        via \texttt{qiskit.quantum\_info.Statevector.evolve}.
      \item Marginalisation over the value register $F$, producing the joint
        $\Pr[(a,b)]$ table.
      \item Paper's Theorem-3 algorithm run with $m'$ oracle queries, plus a
        pure $b{=}0$-rejection variant (see \S\ref{sec:paper-issue}).
      \item Ettinger--H{\o}yer estimator
        $\tilde k = \arg\max_{k \in [1..N/2]} \sum_i \cos(2\pi k a_i / N)$.
      \item Sweep over $m'$ and Monte Carlo repetition
        (400--800 trials per $m'$) to estimate the empirical success curve.
      \item Uniform-$\mathbb{Z}_N$ classical baseline with the same
        post-processing.
    \end{enumerate}
  \item Invocations:
\begin{lstlisting}
python report/evidence/ettinger_hoyer_dihedral.py --N 4  --k0 1  --trials 800 --out report/evidence/results_N4_k1.json
python report/evidence/ettinger_hoyer_dihedral.py --N 8  --k0 3  --trials 800 --out report/evidence/results_N8_k3.json
python report/evidence/ettinger_hoyer_dihedral.py --N 16 --k0 5  --trials 500 --out report/evidence/results_N16_k5.json
python report/evidence/ettinger_hoyer_dihedral.py --N 32 --k0 11 --trials 300 --out report/evidence/results_N32_k11.json
\end{lstlisting}
  \item Verified against Lemma~4: total-variation distance between simulated
        joint and analytical formula.
  \item Wrote up findings and open questions.
\end{enumerate}

\section{Results vs paper}
\subsection{C1 --- Lemma 4 measurement distribution}
\begin{center}
\begin{tabular}{@{}rrrl@{}}
\toprule
$N$ & $k_0$ & TV distance sim vs Lemma~4 & Verdict \\
\midrule
4  & 1  & $3.1\times10^{-16}$ & Match (machine precision) \\
8  & 3  & $4.6\times10^{-16}$ & Match \\
16 & 5  & $5.8\times10^{-16}$ & Match \\
32 & 11 & $6.6\times10^{-15}$ & Match \\
\bottomrule
\end{tabular}
\end{center}
The joint distribution over the first two registers, after full statevector
simulation of $V^\gamma|0\rangle|0\rangle|0\rangle$ and tracing out the value
register, matches the paper's analytical formula to within floating-point
round-off across all sizes tested. Per-$a$ marginals are recorded in
\texttt{report/evidence/results\_N*.json} under the \texttt{per\_a\_marginals}
key.

\subsection{C2 --- Query complexity of $k_0$ recovery}
\begin{center}
\begin{tabular}{@{}rrrrrrr@{}}
\toprule
$N$ & $k_0$ & $m^{\star}_{2/3}$ & $m^{\star}_{1-1/(2N)}$ & Paper $\lceil 64 \ln N\rceil$ & $89\log_2 N + 7$ & Success at paper $m^{\star}$ \\
\midrule
4  & 1  & 2  & 4  & 89  & 185 & 1.00 \\
8  & 3  & 8  & 30 & 134 & 274 & 1.00 \\
16 & 5  & 11 & 45 & 178 & 363 & 1.00 \\
32 & 11 & 18 & 60 & 222 & 452 & 1.00 \\
\bottomrule
\end{tabular}
\end{center}
$m^{\star}_{2/3}$ = smallest $m'$ (oracle queries) at which empirical
success $\geq 2/3$; $m^{\star}_{1-1/(2N)}$ = smallest $m'$ at which empirical
success $\geq 1 - 1/(2N)$ (paper's bound). Observed $m^{\star}$ grows
sub-logarithmically in $N$ over the swept sizes and lies well below both the
Hoeffding-bound-derived $\lceil 64\ln N\rceil$ and the total-cost bound
$89\log N + 7$. This is consistent with the paper (Hoeffding gives a loose
upper bound; empirical constants are much smaller). At the paper-specified
budget, empirical success is $1.00$ in all cases, confirming the promised
$1-1/(2N)$ threshold with margin.

\subsection{C3 --- Quantum vs classical separation}
Uniform baseline (draw $m$ elements of $\mathbb{Z}_N$ i.i.d.\ uniform, apply
the same $\arg\max$ post-processing) reaches success $\geq 2/3$ only at $N=4$
(trivial regime, $m^{\star}_{2/3}=1$ because half of the search space is
already correct); for $N \geq 8$ it \emph{never} reaches $2/3$ within the swept
budget (up to $\sim 8N$ samples). Full per-$m'$ curves are stored in
\texttt{results\_N*.json} under \texttt{classical\_baseline\_success\_prob}.
This is a small-$N$ demonstration of the quantum-vs-classical query gap the
paper's Theorem~3 asserts (the paper cites Simon-style lower bounds; we do
not prove them, only observe the empirical gap on the exact same instances).

\section{The paper's $b{=}1$/sin branch: an inconsistency}\label{sec:paper-issue}
The proof of Theorem~3 says (page~6, our line~262 of the pdftotext extract):
\begin{quote}
``If $m < m'/2$, then the algorithm performs the same classical
post-processing, except that it uses the $m' - m$ measurements for which the
output in the second register is 1, and except that it now seeks to maximize
the sum $\sum_{i=1}^m \sin(2\pi \tilde k a_i / N)$.''
\end{quote}
Under $\Pr[a \mid b{=}1] = \tfrac{2}{N} \sin^2(\pi k_0 a / N)$, symmetry
$a \leftrightarrow N - a$ makes the distribution symmetric under negation,
and $\sin(2\pi k \cdot)$ is an odd function, so
$\mathbb{E}[\sin(2\pi k a / N) \mid b{=}1] = 0$ for \emph{every} $k$.
We verified this numerically (see \texttt{report/evidence/} probe outputs):
for $N{=}16, k_0{=}5$, the $b{=}1$/sin branch achieves only $\sim 10\%$
success at $m=60$, versus $100\%$ for the $b{=}0$/cos branch. The
sin-based estimator cannot discriminate $k_0$.

Our replication mitigates this by running the paper's algorithm literally
\emph{and} a $b{=}0$-only rejection variant that only counts $b{=}0$ shots.
By Hoeffding concentration, $\Pr[m < m'/2]$ decays exponentially in $m'$
(since $\Pr[b{=}0] = 1/2$ exactly when $2k_0 \neq N$), so the paper's overall
success bound $1 - 1/(2N)$ still goes through if one simply \emph{restarts}
the algorithm when the $b{=}0$ branch is unavailable, or if one uses the
squared-cos-in-post-processing trick from the follow-up Ettinger--H{\o}yer /
Regev / Kuperberg literature. The paper's headline query-complexity result is
therefore not undermined, but the specific $b{=}1$/sin construction in the
Theorem-3 proof cannot be taken at face value. We flag this as an interesting
correctable step rather than a fatal flaw.

\section{Verdict and justification}
\textbf{Verdict: REPLICATED.}
\begin{itemize}
  \item C1 (Lemma 4): matched to machine precision on 4 sizes.
  \item C2 (Theorem 3 query bound): empirically confirmed at $100\%$ success
        with query budgets well inside the paper's $O(\log N)$ ceiling.
  \item C3 (quantum vs classical): observed gap on identical instances.
  \item C4, C5: not tested (auxiliary/open).
\end{itemize}
The single caveat --- the $b{=}1$/sin branch of the Theorem-3 proof --- is
noted but does not invalidate the headline results, which rely on the $b{=}0$
branch that concentrates whenever the total oracle budget $m'$ is $\Omega(1)$.
Verdict downgrade candidates: none. This is not SPOT-CHECK (we ran the full
$V^\gamma$ circuit, not just a demo); not PARTIAL (both testable headline
claims C1 and C2 reproduced).

\section{Open Questions}
See \texttt{report/open\_questions.json} for machine-readable form.

\begin{itemize}
  \item[Q1] \textbf{Corrected $b{=}1$ post-processing.} Under
    $\Pr[a|b{=}1] = (2/N)\sin^2(\pi k_0 a / N)$, what is a \emph{provably
    consistent} estimator with success bounds matching the $b{=}0$ branch?
    Candidates: $\arg\max_k |\sum_i \sin(2\pi k a_i/N)|^2$, or reduction to
    the $b{=}0$ distribution via $\sin^2 = 1 - \cos^2$. \emph{Basis:} the
    paper's stated sin-argmax has expectation zero everywhere; the algorithm's
    written form is not consistent as stated even though the overall bound
    still holds by restart/branch-choice arguments. \emph{Next steps:}
    (i) formalize $\arg\max_k |\sum_i \sin(2\pi k a_i/N)|^2$ as a
    hypothesis test with an explicit Hoeffding bound; (ii) empirically
    sweep it against $b{=}0$/cos on the same joints in our harness.
  \item[Q2] \textbf{Empirical constant vs Hoeffding.} We observe
    $m^{\star}_{1-1/(2N)}$ growing much more slowly than
    $\lceil 64\ln N\rceil$. What is the true optimal constant $C$ in
    $m \geq C\ln N$? Fit a curve to the sweep, then run larger $N$
    (via a faster oracle representation) to test whether
    $m^{\star} \sim c_1 \ln N + c_2$ with $c_1 \ll 64$.
    \emph{Basis:} empirical constants are $\sim 3$--$10$ across
    $N\in\{4,\dots,32\}$ vs the paper's $64$; the Hoeffding bound is
    plainly loose. \emph{Next steps:} implement oracle via a permutation
    of computational basis states (no explicit $2^{2n+1}\times 2^{2n+1}$
    matrix), scale to $N=64,128,256$, and fit.
  \item[Q3] \textbf{Post-processing efficiency.} The paper leaves open
    whether classical post-processing can be done in $\mathrm{poly}(\log N)$.
    Kuperberg's 2003 subexponential sieve improved this to
    $2^{O(\sqrt{\log N})}$; Regev linked it to lattice problems. On the
    small $N$ we simulate, the $O(N)$ argmax is trivially fast; is there
    an $O(\mathrm{polylog}(N))$ \emph{approximation} that recovers the
    correct $\tilde k$ with high probability at least for structured $k_0$
    (e.g. small hamming weight)? \emph{Basis:} the paper explicitly flags
    this gap; our sweep confirms the estimator is well-defined but scales
    with $N$. \emph{Next steps:} implement Kuperberg's sieve on $N=64$
    and compare oracle count + wall time to our exhaustive argmax.
  \item[Q4] \textbf{Sensitivity to $\gamma$ choice.} We chose one specific
    left-coset labelling for $\gamma$. Does the concrete choice of
    coset-representative encoding change the observed constants? The
    theoretical distribution is choice-independent (proved), but numerical
    conditioning of $F_N$ and oracle unitarity could differ. \emph{Basis:}
    only one $\gamma$ was tested per $N$; the paper is silent on this. \emph{Next
    steps:} sweep several $\gamma$s per instance (e.g. random-permutation
    coset labels) and compare TV distance + success curves.
  \item[Q5] \textbf{Non-power-of-two $N$.} The paper is stated for general
    $N$, but our implementation is cleanest for $N$ a power of two (uses an
    $n$-qubit $\mathbb{Z}_N$ register). What is the query-complexity constant
    for prime $N$ (which have fewer subgroups)? A specialized qudit or
    exact-DFT-with-padding implementation would let us test whether the
    Hoeffding constant $64$ tightens for prime $N$. \emph{Basis:} our sim
    covers only powers of two, but the paper's promise is general-$N$.
    \emph{Next steps:} implement a qudit or block-encoded $F_N$ for $N\in\{3,5,7,11\}$
    and rerun the sweep, compare $m^{\star}$ scaling by prime vs
    highly composite $N$.
\end{itemize}

\section*{Data and code}
All raw outputs live in \texttt{report/evidence/results\_N*\_k*.json}.
The main script is \texttt{report/evidence/ettinger\_hoyer\_dihedral.py}
(single-file, \textasciitilde{}250 lines, Qiskit 2.5.0). PDF at \texttt{paper.pdf},
pdftotext at \texttt{work/paper.txt}. See \texttt{report/artifacts\_summary.md}
for the full inventory.

\end{document}
